%-------------------------
% Resume in Latex
% Based off of: https://github.com/jakeryang/resume
%------------------------

\documentclass[letterpaper,11pt]{article}

\usepackage{latexsym}
\usepackage[empty]{fullpage}
\usepackage{titlesec}
\usepackage{marvosym}
\usepackage[usenames,dvipsnames]{color}
\usepackage{verbatim}
\usepackage{enumitem}
\usepackage[hidelinks]{hyperref}
\usepackage{fancyhdr}
\usepackage[english]{babel}
\usepackage{tabularx}

% fontawesome
\usepackage{fontawesome5}

% fixed width
\usepackage[scale=0.90,lf]{FiraMono}

% light-grey
\definecolor{light-grey}{gray}{0.83}
\definecolor{dark-grey}{gray}{0.3}
\definecolor{text-grey}{gray}{.08}

\DeclareRobustCommand{\ebseries}{\fontseries{eb}\selectfont}
\DeclareTextFontCommand{\texteb}{\ebseries}

% custom underline
\usepackage{contour}
\usepackage[normalem]{ulem}
\renewcommand{\ULdepth}{1.8pt}
\contourlength{0.8pt}
\newcommand{\myuline}[1]{%
  \uline{\phantom{#1}}%
  \llap{\contour{white}{#1}}%
}

% custom font: helvetica-style
\usepackage{tgheros}
\renewcommand*\familydefault{\sfdefault} 
\usepackage[T1]{fontenc}

\pagestyle{fancy}
\fancyhf{}
\fancyfoot{}
\renewcommand{\headrulewidth}{0pt}
\renewcommand{\footrulewidth}{0pt}

% Adjust margins
\addtolength{\oddsidemargin}{-0.5in}
\addtolength{\evensidemargin}{0in}
\addtolength{\textwidth}{1in}
\addtolength{\topmargin}{-.5in}
\addtolength{\textheight}{1.0in}
\setlength{\footskip}{4.1pt}

\urlstyle{same}

\raggedbottom
\raggedright
\setlength{\tabcolsep}{0in}

% sans serif sections
\titleformat {\section}{
    \bfseries \vspace{2pt} \raggedright \large
}{}{0em}{}[\color{light-grey} {\titlerule[2pt]} \vspace{-4pt}]

%-------------------------
% Custom commands
\newcommand{\resumeItem}[1]{
  \item\small{
    {#1 \vspace{-1pt}}
  }
}

\newcommand{\resumeSubheading}[4]{
  \vspace{-1pt}\item
    \begin{tabular*}{\textwidth}[t]{l@{\extracolsep{\fill}}r}
      \textbf{#1} & {\color{dark-grey}\small #2}\vspace{1pt}\\
      \textit{#3} & {\color{dark-grey} \small #4}\\
    \end{tabular*}\vspace{-4pt}
}

\newcommand{\resumeProjectHeading}[2]{
    \item
    \begin{tabular*}{\textwidth}{l@{\extracolsep{\fill}}r}
      #1 & {\color{dark-grey}\small #2} \\
    \end{tabular*}\vspace{-4pt}
}

% Publication heading: title first, then venue/date metadata.
\newcommand{\resumePubHeading}[3]{
    \item
    \begin{tabularx}{\textwidth}[t]{@{}X@{}}
      \textbf{#1} \\
    \end{tabularx}\vspace{-3pt}
    \begin{tabular*}{\textwidth}{l@{\extracolsep{\fill}}r}
      {\color{dark-grey}\small\textit{#2}} & {\color{dark-grey}\small #3} \\
    \end{tabular*}\vspace{-5pt}
}

% Hardware heading: name + type, no date
\newcommand{\resumeHardwareHeading}[2]{
    \item
    \begin{tabular*}{\textwidth}{l@{\extracolsep{\fill}}r}
      \textbf{#1} --- \textit{#2} & \\
    \end{tabular*}\vspace{-4pt}
}

% Keywords line
\newcommand{\resumeKeywords}[1]{
    \vspace{-2pt}\hspace{18pt}{\small\color{dark-grey}\textit{Keywords: #1}}\vspace{2pt}
}

\newcommand{\resumeSubItem}[1]{\resumeItem{#1}\vspace{-4pt}}

\renewcommand\labelitemii{$\vcenter{\hbox{\tiny$\bullet$}}$}

\newcommand{\resumeSubHeadingListStart}{\begin{itemize}[leftmargin=0in, label={}]}
\newcommand{\resumeSubHeadingListEnd}{\end{itemize}}
\newcommand{\resumeItemListStart}{\begin{itemize}[leftmargin=0.16in, itemsep=0pt, parsep=0pt, topsep=1pt]}
\newcommand{\resumeItemListEnd}{\end{itemize}\vspace{-1pt}}

\color{text-grey}

%-------------------------------------------
%%%%%%  RESUME STARTS HERE  %%%%%%%%%%%%%%%%%%%%%%%%%%%%

\begin{document}

%----------HEADING----------
\begin{center}
    \textbf{\Huge Yu Xia} \\ \vspace{5pt}
    \small \faEnvelope \hspace{2pt} \texttt{yxia072@ucr.edu} \hspace{1pt} $|$
    \hspace{1pt} \faMapMarker* \hspace{2pt}\texttt{Riverside, CA} \hspace{1pt} $|$
    \hspace{1pt} \faLinkedin \hspace{2pt}\texttt{linkedin.com/in/yu-xia-b34074262}
    \\ \vspace{-3pt}
\end{center}


%-----------EDUCATION-----------
\section{EDUCATION}
  \resumeSubHeadingListStart
    \resumeSubheading
      {University of California, Riverside}{Sep. 2023 -- Jun. 2025}
      {Master of Science in Robotics}{Riverside, CA, USA}
    \resumeSubheading
      {Wuhan Huaxia University of Technology}{Sep. 2017 -- Jun. 2021}
      {Bachelor of Engineering in Computer Science and Technology}{Wuhan, China}
  \resumeSubHeadingListEnd

%-----------PUBLICATIONS-----------
\section{PUBLICATIONS}
  \resumeSubHeadingListStart

    \resumePubHeading
      {BOLT: Bandwidth-Efficient Object-Level Streaming for XR Robot Teleoperation}{SenSys 2027 (Under Review)}{Jun. 2026}
      \resumeItemListStart
        \resumeItem{\textbf{Yu Xia}, Hang Qiu}
        \resumeItem{Designed an object-level XR streaming framework for robot teleoperation that transmits object 6D poses and dynamically generated meshes instead of full 3D scenes under strict motion-to-photon latency budgets.}
        \resumeItem{Built a full-stack Meta Quest prototype, reducing uplink traffic to \textbf{0.4 Mbps} ($\sim$220$\times$ lower than point cloud streaming) and improving noisy-pose translation error from \textbf{67.0} to \textbf{25.1 mm} via cross-view refinement.}
      \resumeItemListEnd

    \resumePubHeading
      {Towards Generalizable Safety in Crowd Navigation via Conformal Uncertainty Handling}{CoRL 2025}{Sep. 2024}
      \resumeItemListStart
        \resumeItem{Jianpeng Yao, Xiaopan Zhang*, \textbf{Yu Xia*}, Zejin Wang, Amit K. Roy-Chowdhury, Jiachen Li}
        \resumeItem{A conformal inference--guided constrained RL framework for crowd navigation, achieving \textbf{96.93\% success rate} and \textbf{3.72$\times$ fewer collisions} with robust out-of-distribution performance.}
        \resumeItem{Led the real-robot deployment: built the ROSMASTER X3 sim-to-real stack with LiDAR-based pedestrian detection feeding real-time 3D positions to the RL policy, and conducted crowd-interaction experiments in both sparse and dense environments.}
      \resumeItemListEnd

    \resumePubHeading
      {NL2SpaTiaL: Generating Geometric Spatio-Temporal Logic Specifications from Natural Language for Manipulation Tasks}{IROS 2026}{Nov. 2025}
      \resumeItemListStart
        \resumeItem{Licheng Luo, Kaier Liang, \textbf{Yu Xia}, Mingyu Cai}
        \resumeItem{A hierarchical translation framework that converts natural language into verifiable geometric spatio-temporal logic (SpaTiaL) specifications for manipulation tasks.}
        \resumeItem{Conducted early-stage real-robot validation on ALOHA 2, evaluating $\pi_0$-generated manipulation trajectories against translated SpaTiaL specifications to verify spatio-temporal constraint satisfaction.}
      \resumeItemListEnd

  \resumeSubHeadingListEnd

%-----------ROBOT HARDWARE EXPERIENCE-----------
\section{ROBOT HARDWARE EXPERIENCE}
  \resumeSubHeadingListStart

    \resumeHardwareHeading{TurtleBot 4}{Differential-Drive Mobile Robot}
      \resumeItemListStart
        \resumeItem{Differential-drive mobile robot with iRobot Create 3 base, Raspberry Pi 4, OAK-D stereo camera, and RPLIDAR A1M8 2D LiDAR.}
        \resumeItem{Used in earlier real-world experiments of the crowd navigation work (CoRL 2025), running the same LiDAR-based pedestrian detection pipeline for real-time avoidance in dynamic indoor environments.}
      \resumeItemListEnd

    \resumeHardwareHeading{Yahboom ROSMASTER X3}{Omnidirectional Mobile Robot}
      \resumeItemListStart
        \resumeItem{Mecanum-wheel omnidirectional mobile robot with Jetson Orin, SLAMTEC LiDAR, and Astra depth camera.}
        \resumeItem{Primary sim-to-real platform for the CoRL 2025 crowd navigation deployment: LiDAR-based pedestrian detection provides real-time 3D pedestrian positions to the RL policy for avoidance in both sparse and dense crowds.}
      \resumeItemListEnd

    \resumeHardwareHeading{ALOHA 2}{Bimanual Teleoperation System}
      \resumeItemListStart
        \resumeItem{Low-cost bimanual teleoperation system with 4$\times$ Trossen ViperX/WidowX 6-DoF arms, leader-follower control, and multi-view RGB cameras.}
        \resumeItem{Used for early-stage real-robot evaluation in the NL2SpaTiaL project (IROS 2026): action trajectories generated by $\pi_0$ are evaluated against translated SpaTiaL formulas to verify whether executed manipulations satisfy the spatio-temporal specifications.}
      \resumeItemListEnd

    \resumeHardwareHeading{Kinova Gen3}{7-DoF Collaborative Arm}
      \resumeItemListStart
        \resumeItem{7-DoF ultra-lightweight collaborative arm (8.2\,kg, 902\,mm reach, 4\,kg payload) with integrated torque sensors and 2D/3D vision module.}
        \resumeItem{Serving as the manipulation execution end for an ongoing VR teleoperation project. Integrating Kortex API with WebRTC-based remote control for real-time manipulation over WAN.}
      \resumeItemListEnd

    \resumeHardwareHeading{Hexapod Robot (DIY)}{Legged Locomotion Platform}
      \resumeItemListStart
        \resumeItem{Custom-built six-legged robot for RL-based locomotion research, including gait training and flip-recovery behavior.}
        \resumeItem{Trained RL locomotion and fall-recovery policies in simulation. Developed the real-robot controller for deployment; sim-to-real transfer pending due to hardware actuator limitations.}
      \resumeItemListEnd

  \resumeSubHeadingListEnd

%-----------PROJECT EXPERIENCE-----------
\section{PROJECT EXPERIENCE}
  \resumeSubHeadingListStart

    \resumeProjectHeading
      {\textbf{Robotics Fabrication: Ball Balancing Car}}{Apr. 2024}
      \resumeItemListStart
        \resumeItem{\textbf{Role}: Perception \& Control Developer. Built subsystems for sensing, wireless communication, and PID motor control for real-time ball balancing.}
      \resumeItemListEnd
      \resumeKeywords{Arduino, PID, IMU, Bluetooth}

    \resumeProjectHeading
      {\textbf{UAV Path Planning and Obstacle Avoidance}}{Dec. 2023}
      \resumeItemListStart
        \resumeItem{\textbf{Role}: Simulator \& Communication Developer. Set up simulator, configured UAV sensors, established ROS--PX4 communication for RL training.}
        \resumeItem{Developed a rule-based baseline for navigating obstacles; contributed to RL reward function design to accelerate convergence.}
      \resumeItemListEnd
      \resumeKeywords{ROS, PX4, Gazebo, PPO}

    \resumeProjectHeading
      {\textbf{Mask Wearing Detection with YOLOv5}}{Sep. 2021}
      \resumeItemListStart
        \resumeItem{Trained and optimized YOLOv5 for face mask verification, achieving \textbf{85\% accuracy} on SF-MASK via iterative hyperparameter tuning.}
      \resumeItemListEnd
      \resumeKeywords{YOLOv5, PyTorch, OpenCV}

    \resumeProjectHeading
      {\textbf{Image Classification System Design}}{May 2021}
      \resumeItemListStart
        \resumeItem{Enhanced AlexNet architecture and achieved \textbf{85\% accuracy} on CIFAR-10 through comparative experiments using TensorFlow.}
      \resumeItemListEnd
      \resumeKeywords{TensorFlow, AlexNet, CIFAR-10}

  \resumeSubHeadingListEnd

%-----------SKILLS-----------
\section{SKILLS}
 \begin{itemize}[leftmargin=0in, label={}]
    \small{\item{
     \textbf{Programming}{: Python, C/C++, Java, MATLAB, HTML5, SQL}\vspace{2pt} \\
     \textbf{Robotics \& Sim}{: ROS / ROS\,2, MoveIt\,2, Gazebo, MuJoCo, Isaac Sim, PX4, Unity}\vspace{2pt} \\
     \textbf{AI / ML}{: PyTorch, TensorFlow, OpenCV, MediaPipe, YOLO, Reinforcement Learning}\vspace{2pt} \\
     \textbf{Tools}{: Docker, Git, Linux (Ubuntu), WebRTC, SolidWorks, Blender, Photoshop}\vspace{2pt} \\
     \textbf{Robot Platforms}{: TurtleBot 4, ROSMASTER X3, ALOHA 2, Kinova Gen3, Hexapod (DIY), UAV (PX4)}
    }}
 \end{itemize}

%-----------SUMMARY-----------
\section{SUMMARY}
\small{Robotics M.S. from UC Riverside with work at CoRL 2025 and IROS 2026, and BOLT under review at SenSys 2027. Hands-on experience in robot manipulation, deep reinforcement learning, XR teleoperation, and sim-to-real deployment across 6 physical robot platforms. Proficient in Python, C++, ROS\,2, MoveIt\,2, PyTorch, and simulation tools (Gazebo, MuJoCo, Isaac Sim). U.S. Permanent Resident.}
%-------------------------------------------
\end{document}
